\section{CONCLUSION}
This paper reformulated collaborative CPS recovery as a trustworthy heterogeneous federated-recovery problem in which both physical observations and federated recovery knowledge may be untrusted. The revised design separates three responsibilities that were previously mixed together: heterogeneous local adaptation, Byzantine-resilient knowledge coordination, and physical execution safety. Subsystem-specific observation and action adapters remain local, only a shared recovery backbone is federated, and representation alignment is used to reduce honest-client drift rather than authenticate clients.

The security mechanism no longer relies on client-reported credibility as a trust root. Honest clients gate recovery knowledge with local utility and safety-recovery certificates, while the server applies Byzantine-resilient aggregation under an explicit malicious-client budget. For communication-efficient Top-$k$ sharing, the compression residual is characterized by the deterministic contraction factor $\delta_k=k/d$ instead of being introduced as an unverifiable assumption. The resulting optimization bound therefore keeps both compression and Byzantine residual terms explicit.

Most importantly, the learned federated action is treated as an untrusted proposal rather than a directly executable command. A robust safety-recovery shield enforces the physical contract before actuation. Under uncertainty-set coverage and shield feasibility, the safe set is forward invariant for arbitrary learned policy proposals, including actions produced by a poisoned global model. This separates physical safety from successful convergence of the learning algorithm and gives the recovery framework a mechanism-level guarantee that empirical robust aggregation alone cannot provide.

The existing IEEE 39-bus and IEEE 118-bus results remain evidence for recovery performance, communication efficiency, scalability, heterogeneity stress, and empirical robustness under the attacks actually implemented. They are not reinterpreted as proof of the newly introduced certificate guarantees. The revised security-validation protocol therefore requires explicit barrier-margin logging, unsafe-policy injection, adaptive and colluding Byzantine attacks, aggregation-deviation auditing, code-to-claim traceability, and privacy-budget reporting when the optional differential-privacy extension is enabled. These additions define the path from a communication-efficient heterogeneous FRL prototype to a high-assurance trustworthy recovery mechanism for safety-critical CPS.
